\chapter{Von der ebenen Pose zum räumlichen Eisenbahnzustand}
\label{chap:pose3-and-cant}

An der gewählten Übergangsbogenposition liefert pose2 ebene Lage und Richtung.
Ein Eisenbahnquerschnitt kann damit noch nicht angeordnet werden: Das
horizontale Krümmungsgesetz des Alignments enthält weder Höhe noch Drehung der
Schienenebene. Die räumliche Interpretation benötigt deshalb zusätzliche,
eigenständig qualifizierte Längsgesetze.

\section{Zusätzliche Gesetze, nicht zusätzliche Identität}

\(h(s)\) bezeichne das Höhengesetz in der gewählten metrischen Realisierung.
Soweit unterstützt, liefert \(h'(s)\) die Neigung und \(h''(s)\) unter der
verwendeten Profilkonvention einen Beitrag zur vertikalen
Krümmungsinformation. \(u(s)\) bezeichne die Überhöhung als
Schienenhöhendifferenz:

\begin{definition}[Überhöhungshöhendifferenz]
\label{def:cant-height-difference}
\[
u\colon D_u\rightarrow \mathbb R
\]
ist ein Längsgesetz der Überhöhung mit eigenem Definitionsbereich \(D_u\),
Einheiten, eigener Partition, Provenienz und Anwendbarkeit.
\end{definition}

Überhöhung ist keine horizontale Krümmung. Die Gesetze \(\kappa(s)\), \(h(s)\)
und \(u(s)\) können unterschiedliche Abschnittsgrenzen und Quellen besitzen.
Ihr gemeinsamer Parameter ermöglicht die synchrone Auswertung; er macht sie
weder zu einer Eigenschaft noch verleiht er einem Gesetz Autorität über ein
anderes.

\begin{definition}[Überhöhungsreferenzkonvention]
\label{def:cant-reference-convention}
Eine Überhöhungskonvention gibt Vorzeichen, Drehachse,
Spurweiteninterpretation, Bezugslinie oder Bezugsschiene sowie eine mögliche
Verschiebung der Bezugslinie an.
\end{definition}

Nur für die hier verwendete Erklärungskonvention wird \(u\) über eine
angegebene Spurweite \(g_{\mathrm{gauge}}>0\) gemessen, die Mittellinie bleibt
Bezug, und positiver Roll folgt dem dargestellten rechtshändigen Frame. Dann
gilt
\begin{equation}
\phi(s)=
\arctan\!\left(\frac{u(s)}{g_{\mathrm{gauge}}}\right).
\label{eq:derived-roll-angle}
\end{equation}
Diese Gleichung wandelt eine qualifizierte Darstellung der Überhöhung in einen
Rollwinkel um. Sie ist keine universelle Überhöhungskonvention und bestimmt
nicht, welches Überhöhungsgesetz zu entwerfen ist.

\begin{figure}[htbp]
\centering
\input{figures/state/cant_interpretation-DE}
\caption{Lokaler Querschnitt: Überhöhung dreht das Quer--Vertikal-Paar um die Tangente; Konvention und Spurweite bleiben explizite Eingaben.}
\label{fig:cant-interpretation}
\end{figure}

\section{Das bewegte Eisenbahn-Frame}

\begin{definition}[Nicht überhöhtes geometrisches Frame]
\label{def:uncanted-frame}
An \(s\) sei
\[
F_0(s)=\bigl(\mathbf t(s),\mathbf l_0(s),\mathbf z_0(s)\bigr)
\]
ein orientiertes orthonormales Frame: \(\mathbf t\) ist longitudinal,
\(\mathbf l_0\) und \(\mathbf z_0\) sind Quer- und Vertikalrichtung aus einer
angegebenen Realisierungs- und Transportkonvention.
\end{definition}

Horizontale Krümmung steuert die Richtungsänderung. Vertikale Geometrie steuert
die Neigung von \(\mathbf t\). Die metrische Realisierung liefert räumlichen
Bezug und Vertikalinterpretation. Eine Stetigkeitsregel transportiert den
Querschnitt durch Geraden und Bereiche kleiner Krümmung, in denen eine reine
Frenet-Konstruktion instabil wäre.

\begin{definition}[Eisenbahn-Frame]
\label{def:railway-frame}
Das überhöhungsbewusste Gleis-Frame lautet
\[
F_{\mathrm{rail}}(s)=
\bigl(\mathbf t(s),\mathbf l_\phi(s),\mathbf z_\phi(s)\bigr),
\qquad
(\mathbf l_\phi,\mathbf z_\phi)
=R_{\mathbf t,\phi(s)}(\mathbf l_0,\mathbf z_0).
\]
\end{definition}

Diese lokale rechtshändige Konstruktion wird gewählt, um Abhängigkeiten
sichtbar zu machen; sie begründet keine universelle AIM-Frame-Konvention.
Quellkonventionen müssen vor einem Vergleich von Vorzeichen oder Schienenlagen
explizit transformiert werden.

\begin{definition}[Fahrzeug-Frame]
\label{def:vehicle-frame}
\(F_{\mathrm{veh}}(s)\) bezeichnet ein optionales fahrzeugabhängiges Frame,
das aus Gleis-Frame und angegebenem Fahrzeug- und Federungsmodell erzeugt wird.
\end{definition}

Ein Fahrzeug-Frame lässt sich nicht allein aus Alignment-Geometrie ableiten.

\section{Was pose3 ergänzt}

\begin{definition}[Laufzeit-pose3]
\label{def:runtime-pose3}
Für diese Erklärung auf Thesis-Ebene lautet die räumliche Eisenbahnpose an
\(s\)
\begin{equation}
\mathrm{pose3}(s)=
\bigl(\Gamma(s),F_{\mathrm{rail}}(s),
h'(s),h''(s),u(s)\bigr),
\label{eq:runtime-pose3}
\end{equation}
wobei nicht verfügbare oder nicht anwendbare Profilglieder entfallen und alle
Größen durch ihre Realisierung und Konvention qualifiziert sind.
\end{definition}

\(\Gamma(s)\) ist die aus pose2, Profil, metrischer Interpretation und
Koordinatendiensten zusammengesetzte räumliche Referenztrajektorie. Der
Frame-Operator \(\mathcal F\) aus \cref{def:frame-operator} beschreibt die
typisierte Auswertung. pose3 ergänzt die zur Orientierung eines
Eisenbahnquerschnitts benötigten Größen; es ist weder die vollständige
physische Anlage noch umfasst es Material, Zustand, Topologie,
Fahrzeugverhalten oder Engineering-Autorität.

\begin{proposition}[pose3-Rekonstruktionsprinzip]
\label{prop:pose3-reconstruction}
Aus \(\kappa(s)\), \(h(s)\) und \(u(s)\) sowie Anfangspose, metrischer
Realisierung, Spurweite, Orientierungs- und Frame-Konventionen können die
anwendbaren pose3-Felder ausgewertet werden. Das Ergebnis bewahrt seinen
Alignment-Bezug und die Provenienz jeder zusätzlichen Eingabe.
\end{proposition}

\section{Soll, Realisierung, physischer Zustand und Beobachtung}

Das beabsichtigte Alignment und seine Entwurfsgesetze erzeugen eine
\emph{Soll-Pose}. Eine Metric Realization ordnet dieses Soll in einem
metrischen Raum an. Die Ausführung erzeugt eine physische Anlage, deren Zustand
sich mit der Zeit ändert. Vermessung erzeugt zeitgestempelte Beobachtungen oder
Schätzungen des beobachteten Zustands in dargestellten Koordinaten. Diese
Größen sind verbunden, aber nicht austauschbar.

Eine CAD-Kurve kann die Solltrajektorie darstellen. Eine
Vermessungspolylinie kann Beobachtungen der Schienen darstellen. Ein
Soll--Ist-Vergleich muss beide Seiten durch Objektbezug, Zeit, Realisierung,
Konvention, Provenienz und Unsicherheit qualifizieren, bevor eine Abweichung
interpretiert werden kann. Koordinatensubtraktion allein kann weder den
richtigen noch den akzeptablen Zustand bestimmen.

\section{Wo Überhöhung und Krümmung zusammentreffen}

Krümmung und Überhöhung bleiben getrennte Gesetze, auch wenn eine betriebliche
Berechnung beide verwendet. Die Beziehungen in
\cref{eq:effective-lateral-acceleration,eq:cant-equilibrium,eq:jerk-relation}
sind vereinfachte, annahmenabhängige Interpretationen für die folgende
Dynamikdiskussion. Sie sind keine autoritativen Entwurfskopplungsregeln.
Anwendbare Normen, Fahrzeugmodelle, Geschwindigkeitsprofile und Entscheidungen
bleiben getrennte Eingaben.

\begin{summary}
pose3 ist ein qualifizierter räumlicher Zustand, der aus dem identifizierten
Alignment sowie Profil, Überhöhung, metrischer Realisierung und
Frame-Konvention abgeleitet wird. Überhöhung dreht das Eisenbahn-Frame, ändert
aber weder die Identität von \(\kappa(s)\) noch begründet sie
Engineering-Autorität.
\end{summary}

\begin{proposition}[Gleiszustandsprinzip]
pose3 beschreibt eine qualifizierte Eisenbahnpose. Fahrzeugzustände benötigen
zusätzliche Vehicle-Space-Operatoren und Modelle.
\end{proposition}
